% ------------------------------------------------------------
% Geometry Module
% Deep Mathematical Spine for Structural AI Operators
% ------------------------------------------------------------

\section{Connections and Curvature}

\subsection{Affine Connections}

\begin{definition}
Let $\M$ be a smooth manifold.  
An \emph{affine connection} on $\M$ is a bilinear map
\[
\nabla : \Gamma(T\M) \times \Gamma(T\M) \to \Gamma(T\M)
\]
assigning to each pair of vector fields $(X,Y)$ a vector field
$\nabla_X Y$ such that:
\begin{enumerate}
    \item $\nabla_{fX + gY} Z = f \nabla_X Z + g \nabla_Y Z$,
    \item \nabla_X (fY) = X(f) Y + f \nabla_X Y,
\end{enumerate}
for all smooth functions $f,g$ and vector fields $X,Y,Z$.
\end{definition}

\subsection{Torsion}

\begin{definition}
The \emph{torsion tensor} of a connection $\nabla$ is the $(1,2)$-tensor
\[
T(X,Y) = \nabla_X Y - \nabla_Y X - [X,Y].
\]
\end{definition}

\begin{lemma}
A connection is torsion-free if and only if
\[
\nabla_X Y - \nabla_Y X = [X,Y]
\]
for all vector fields $X,Y$.
\end{lemma}

\subsection{Curvature}

\begin{definition}
The \emph{Riemann curvature tensor} of $\nabla$ is the $(1,3)$-tensor
\[
R(X,Y)Z = \nabla_X \nabla_Y Z - \nabla_Y \nabla_X Z - \nabla_{[X,Y]} Z.
\]
\end{definition}

\begin{lemma}
For any connection $\nabla$, the curvature tensor satisfies the
antisymmetry property
\[
R(X,Y)Z = -R(Y,X)Z.
\]
\end{lemma}

\subsection{Parallel Transport}

\begin{definition}
Let $\gamma : [0,1] \to \M$ be a smooth curve.  
A vector field $V(t)$ along $\gamma$ is said to be \emph{parallel} if
\[
\nabla_{\dot{\gamma}(t)} V(t) = 0.
\]
\end{definition}

\begin{theorem}
For any smooth curve $\gamma$ and initial vector $v_0 \in T_{\gamma(0)}\M$,
there exists a unique parallel vector field $V(t)$ along $\gamma$ with
$V(0) = v_0$.
\end{theorem}

\subsection{Holonomy}

\begin{definition}
The \emph{holonomy group} $\Hol(p)$ at a point $p \in \M$ is the set of
linear transformations of $T_p\M$ obtained by parallel transport along
piecewise smooth loops based at $p$.
\end{definition}

\begin{lemma}
If $\M$ is connected, the holonomy group is generated by parallel
transport along loops contained in any sufficiently small neighborhood
of $p$.
\end{lemma}

\section{Geodesics and the Levi-Civita Connection}

\subsection{Metric Compatibility}

\begin{definition}
Let $(\M,g)$ be a Riemannian manifold.  
A connection $\nabla$ is said to be \emph{metric compatible} if
\[
X\big(g(Y,Z)\big)
= g(\nabla_X Y, Z) + g(Y, \nabla_X Z)
\]
for all vector fields $X,Y,Z$.
\end{definition}

\begin{lemma}
Metric compatibility is equivalent to preservation of inner products under
parallel transport: if $V(t)$ and $W(t)$ are parallel along a curve
$\gamma(t)$, then
\[
\frac{d}{dt} g\big(V(t),W(t)\big) = 0.
\]
\end{lemma}

\subsection{Levi-Civita Connection}

\begin{theorem}[Fundamental Theorem of Riemannian Geometry]
On a Riemannian manifold $(\M,g)$, there exists a unique connection
$\nabla$ that is:
\begin{enumerate}
    \item torsion-free, and
    \item metric compatible.
\end{enumerate}
This connection is called the \emph{Levi-Civita connection}.
\end{theorem}

\begin{proof}[Sketch]
Uniqueness follows from the Koszul formula:
\[
2 g(\nabla_X Y, Z)
= X g(Y,Z) + Y g(Z,X) - Z g(X,Y)
+ g([X,Y],Z) - g([Y,Z],X) - g([Z,X],Y).
\]
Existence is shown by verifying that the right-hand side defines a
connection satisfying torsion-free and metric compatibility conditions.
\end{proof}

\subsection{Geodesics}

\begin{definition}
A smooth curve $\gamma : I \to \M$ is a \emph{geodesic} if its velocity
vector is parallel along the curve:
\[
\nabla_{\dot{\gamma}(t)} \dot{\gamma}(t) = 0.
\]
\end{definition}

\begin{lemma}
In local coordinates $\{x^i\}$, the geodesic equation becomes
\[
\ddot{x}^k + \Gamma^k_{ij} \dot{x}^i \dot{x}^j = 0,
\]
where $\Gamma^k_{ij}$ are the Christoffel symbols of the Levi-Civita
connection.
\end{lemma}

\subsection{Exponential Map}

\begin{definition}
Let $p \in \M$.  
The \emph{exponential map} at $p$ is defined by
\[
\exp_p(v) = \gamma_v(1),
\]
where $\gamma_v$ is the unique geodesic with initial conditions
$\gamma_v(0) = p$ and $\dot{\gamma}_v(0) = v$.
\end{definition}

\begin{theorem}
For each $p \in \M$, the exponential map $\exp_p$ is smooth on a
neighborhood of $0 \in T_p\M$ and satisfies
\[
\frac{d}{dt}\exp_p(tv)\bigg|_{t=0} = v.
\]
\end{theorem}

\subsection{Injectivity Radius}

\begin{definition}
The \emph{injectivity radius} at $p$ is
\[
\mathrm{inj}(p)
= \sup \left\{ r > 0 \;\middle|\; \exp_p : B_r(0) \subset T_p\M \to \M
\text{ is a diffeomorphism onto its image} \right\}.
\]
\end{definition}

\begin{lemma}
If $\mathrm{inj}(p) > 0$, then geodesics starting at $p$ minimize distance
for all $t$ such that $\|\dot{\gamma}(0)\| t < \mathrm{inj}(p)$.
\end{lemma}

\section{Curvature Identities}

\subsection{Sectional Curvature}

\begin{definition}
Let $(\M,g)$ be a Riemannian manifold and let $X,Y \in T_p\M$ be
linearly independent.  
The \emph{sectional curvature} of the plane spanned by $X$ and $Y$ is
\[
K(X,Y)
= \frac{g\big(R(X,Y)Y, X\big)}
       {g(X,X)g(Y,Y) - g(X,Y)^2}.
\]
\end{definition}

\begin{lemma}
If $X$ and $Y$ are orthonormal, then
\[
K(X,Y) = g\big(R(X,Y)Y, X\big).
\]
\end{lemma}

\subsection{Ricci Curvature}

\begin{definition}
The \emph{Ricci curvature} is the trace of the Riemann curvature tensor:
\[
\mathrm{Ric}(X,Y)
= \sum_{i=1}^n g\big(R(e_i,X)Y, e_i\big),
\]
where $\{e_i\}$ is any orthonormal basis of $T_p\M$.
\end{definition}

\begin{lemma}
Ricci curvature is independent of the choice of orthonormal basis.
\end{lemma}

\subsection{Scalar Curvature}

\begin{definition}
The \emph{scalar curvature} is the trace of the Ricci tensor:
\[
S = \sum_{i=1}^n \mathrm{Ric}(e_i,e_i).
\]
\end{definition}

\begin{lemma}
Scalar curvature satisfies
\[
S = \sum_{i,j} g\big(R(e_i,e_j)e_j, e_i\big).
\]
\end{lemma}

\subsection{First Bianchi Identity}

\begin{theorem}[First Bianchi Identity]
For any vector fields $X,Y,Z$,
\[
R(X,Y)Z + R(Y,Z)X + R(Z,X)Y = 0.
\]
\end{theorem}

\begin{proof}[Sketch]
The identity follows from the torsion-free property of the Levi-Civita
connection and the Jacobi identity for Lie brackets.
\end{proof}

\subsection{Second Bianchi Identity}

\begin{theorem}[Second Bianchi Identity]
The covariant derivative of the curvature tensor satisfies
\[
(\nabla_X R)(Y,Z) + (\nabla_Y R)(Z,X) + (\nabla_Z R)(X,Y) = 0.
\]
\end{theorem}

\begin{lemma}
The second Bianchi identity implies that the divergence of the Ricci
tensor satisfies
\[
\nabla^i \mathrm{Ric}_{ij}
= \frac{1}{2} \nabla_j S.
\]
\end{lemma}

\subsection{Curvature Bounds}

\begin{definition}
A manifold $(\M,g)$ has \emph{bounded sectional curvature} if there exist
constants $k_1 \le k_2$ such that
\[
k_1 \le K(X,Y) \le k_2
\]
for all $p \in \M$ and all planes spanned by $X,Y \in T_p\M$.
\end{definition}

\begin{theorem}
If sectional curvature is bounded above by $k_2$, then geodesics have no
conjugate points for distances less than $\pi/\sqrt{k_2}$.
\end{theorem}

\begin{lemma}
If sectional curvature is bounded below by $k_1 > 0$, then the injectivity
radius satisfies
\[
\mathrm{inj}(p) \le \frac{\pi}{\sqrt{k_1}}.
\]
\end{lemma}

\section{Holonomy and Drift}

\subsection{Holonomy Decomposition}

\begin{definition}
Let $(\M,g)$ be a Riemannian manifold with Levi-Civita connection $\nabla$.
For a point $p \in \M$, the \emph{holonomy group} $\Hol(p)$ is the set of
linear maps
\[
P_\gamma : T_p\M \to T_p\M
\]
obtained by parallel transport along piecewise smooth loops $\gamma$
based at $p$.
\end{definition}

\begin{theorem}[Holonomy Decomposition]
Let $\Hol(p)$ act on $T_p\M$.  
Then $T_p\M$ decomposes into $\Hol(p)$-invariant subspaces
\[
T_p\M = V_1 \oplus \cdots \oplus V_k,
\]
each corresponding to an irreducible representation of the holonomy group.
\end{theorem}

\begin{lemma}
If $\M$ is connected, the holonomy decomposition is independent of the
choice of base point up to canonical isomorphism.
\end{lemma}

\subsection{Ambrose--Singer Theorem (Statement)}

\begin{theorem}[Ambrose--Singer]
The holonomy algebra $\mathfrak{hol}(p)$ is generated by curvature
operators of the form
\[
R(X,Y)_p,
\]
where $X,Y$ are parallel transported along curves starting at $p$.
\end{theorem}

\begin{remark}
This theorem links curvature directly to holonomy, providing the geometric
foundation for drift operators defined below.
\end{remark}

\subsection{Holonomy-Induced Transport Deviation}

\begin{definition}
Let $\gamma$ be a smooth curve from $p$ to $q$.  
Let $P_\gamma$ denote parallel transport along $\gamma$.  
Let $H_\gamma$ denote hybrid or discrete-induced transport along the same
path.  
The \emph{transport deviation} along $\gamma$ is
\[
D_\gamma = H_\gamma - P_\gamma.
\]
\end{definition}

\begin{lemma}
If $H_\gamma$ coincides with parallel transport for all $\gamma$, then
$D_\gamma = 0$ and the hybrid system is geometrically faithful.
\end{lemma}

\subsection{Holonomy Drift Tensor}

\begin{definition}
Let $(\M,g)$ be a Riemannian manifold with Levi-Civita connection $\nabla$.
Let $H_\gamma$ be a hybrid transition operator acting on tangent vectors
along curves $\gamma$.  
The \emph{holonomy drift tensor} at $p$ is the $(1,2)$-tensor
\[
\mathcal{D}_p(X,Y)
= \lim_{t \to 0}
\frac{H_{\gamma_{X,Y}(t)} - P_{\gamma_{X,Y}(t)}}{t},
\]
where $\gamma_{X,Y}(t)$ is a two-parameter family of infinitesimal loops
generated by $X,Y \in T_p\M$.
\end{definition}

\begin{lemma}
If $H_\gamma$ is smooth and agrees with parallel transport to first order,
then $\mathcal{D}_p$ is well-defined and tensorial.
\end{lemma}

\subsection{Drift–Curvature Relation}

\begin{theorem}
If hybrid transport differs from geometric parallel transport only through
a curvature-dependent perturbation, then
\[
\mathcal{D}_p(X,Y)
= A_p\big(R(X,Y)\big)
\]
for some linear operator $A_p : \mathrm{End}(T_p\M) \to \mathrm{End}(T_p\M)$.
\end{theorem}

\begin{remark}
This relation shows that drift is controlled by curvature, and therefore
holonomy.  
Hybrid operators inherit geometric structure through this tensor.
\end{remark}

\subsection{Drift Invariants}

\begin{definition}
A quantity $\mathcal{I}(p)$ is a \emph{drift invariant} if
\[
\mathcal{I}(p) = \mathcal{I}(q)
\]
whenever $p$ and $q$ are connected by a curve along which
$\mathcal{D}_\gamma = 0$.
\end{definition}

\begin{lemma}
If $\mathcal{D}_p = 0$ on a neighborhood of $p$, then drift invariants
coincide with holonomy invariants.
\end{lemma}

\section{Geometry of Embedded Graphs}

\subsection{Manifold--Graph Embeddings}

\begin{definition}
Let $(\M,g)$ be a Riemannian manifold and let $G = (V,E)$ be a finite
graph.  
A \emph{manifold--graph embedding} is a map
\[
\iota : V \to \M
\]
such that each edge $(u,v) \in E$ corresponds to a curve $\gamma_{uv}$
in $\M$ connecting $\iota(u)$ and $\iota(v)$.
\end{definition}

\begin{lemma}
If each $\gamma_{uv}$ is chosen to be a minimizing geodesic, then
\[
d_\M(\iota(u),\iota(v)) = \ell(\gamma_{uv}),
\]
where $\ell(\gamma_{uv})$ is the geodesic length.
\end{lemma}

\subsection{Distortion Bounds}

\begin{definition}
The \emph{embedding distortion} of $\iota$ is
\[
\mathrm{dist}(\iota)
= \sup_{u \ne v}
\frac{d_\M(\iota(u),\iota(v))}{d_G(u,v)},
\]
where $d_G$ is the graph distance.
\end{definition}

\begin{lemma}
If sectional curvature is bounded above by $k_2$, then for sufficiently
small neighborhoods,
\[
d_\M(\iota(u),\iota(v))
\le d_G(u,v) \left( 1 + O(k_2\, d_G(u,v)^2) \right).
\]
\end{lemma}

\begin{remark}
Curvature controls how faithfully graph distances can be represented
geometrically.
\end{remark}

\subsection{Geodesic Consistency}

\begin{definition}
An embedding $\iota$ is \emph{geodesically consistent} if for every
three vertices $u,v,w$ with $v$ lying on a shortest path between $u$ and
$w$ in $G$, the point $\iota(v)$ lies on a minimizing geodesic between
$\iota(u)$ and $\iota(w)$.
\end{definition}

\begin{theorem}
If $\iota$ has distortion $< 2$ and sectional curvature is nonpositive,
then $\iota$ is geodesically consistent.
\end{theorem}

\begin{proof}[Sketch]
Nonpositive curvature ensures geodesic convexity, and low distortion
forces midpoint preservation along shortest paths.
\end{proof}

\subsection{Hybrid Transport Compatibility}

\begin{definition}
Let $H$ be a hybrid transition operator acting on pairs $(p,q)$ with
$p \in \M$ and $q \in V$.  
The embedding $\iota$ is \emph{hybrid compatible} if
\[
H(\iota(u),u) = (\iota(v),v)
\]
whenever $(u,v) \in E$.
\end{definition}

\begin{lemma}
If $\iota$ is geodesically consistent and hybrid transitions follow
graph edges, then hybrid transport agrees with geometric transport up to
first order.
\end{lemma}

\begin{remark}
This lemma is the geometric foundation for the holonomy drift tensor
defined in Section~4.
\end{remark}

\subsection{Collapse-Surface Geometry}

\begin{definition}
Let $C \subseteq V$ be a collapse surface in the discrete system.  
The \emph{geometric collapse region} is
\[
\mathcal{C}_\M = \iota(C) \subseteq \M.
\]
\end{definition}

\begin{lemma}
If $\iota$ has bounded distortion, then the geometric collapse region
has bounded diameter:
\[
\mathrm{diam}(\mathcal{C}_\M)
\le \mathrm{dist}(\iota) \cdot \mathrm{diam}(C).
\]
\end{lemma}

\begin{theorem}
If sectional curvature is bounded and $\iota$ is hybrid compatible, then
collapse-avoidance operators can be realized as geometric steering away
from $\mathcal{C}_\M$.
\end{theorem}

\begin{remark}
This theorem is the geometric precursor to the collapse-avoidance
operator in the AI-operators module
